\section{AIME II 2002}
        \begin{enumerate}
        	\item (\textit{AIME II 2002}) Given that $x$ and $y$ are both integers between $100$ and $999$, inclusive; $y$ is the number formed by reversing the digits of $x$; and $z=|x-y|$. How many distinct values of $z$ are possible?


 

	\item (\textit{AIME II 2002}) Three \href{https://artofproblemsolving.com/wiki/index.php?title=Vertex}{vertices} of a \href{https://artofproblemsolving.com/wiki/index.php?title=Cube}{cube} are $P=(7,12,10)$, $Q=(8,8,1)$, and $R=(11,3,9)$. What is the \href{https://artofproblemsolving.com/wiki/index.php?title=Surface\_area}{surface area} of the cube?


 

	\item (\textit{AIME II 2002}) It is given that $\log_{6}a + \log_{6}b + \log_{6}c = 6$, where $a$, $b$, and $c$ are \href{https://artofproblemsolving.com/wiki/index.php?title=Positive}{positive} \href{https://artofproblemsolving.com/wiki/index.php?title=Integer}{integers} that form an increasing \href{https://artofproblemsolving.com/wiki/index.php?title=Geometric\_sequence}{geometric sequence} and $b - a$ is the \href{https://artofproblemsolving.com/wiki/index.php?title=Perfect\_square}{square} of an integer. Find $a + b + c$.


 

	\item (\textit{AIME II 2002}) Patio blocks that are hexagons $1$ unit on a side are used to outline a garden by placing the blocks edge to edge with $n$ on each side. The diagram indicates the path of blocks around the garden when $n=5$.


 
                    \begin{figure}[H]
                        \centering
                        \includegraphics[width=0.5\textwidth]{Problem/AIME_II_2002_4_Problem_Figure_1.png}
                    \end{figure}
                    


 If $n=202$, then the area of the garden enclosed by the path, not including the path itself, is $m\left(\sqrt3/2\right)$ square units, where $m$ is a positive integer. Find the remainder when $m$ is divided by $1000$.


 

	\item (\textit{AIME II 2002}) Find the sum of all positive integers $a=2^n3^m$ where $n$ and $m$ are non-negative integers, for which $a^6$ is not a divisor of $6^a$.


 

	\item (\textit{AIME II 2002}) Find the integer that is closest to $1000\sum_{n=3}^{10000}\frac1{n^2-4}$.


 

	\item (\textit{AIME II 2002}) It is known that, for all positive integers $k$,


 $1^2+2^2+3^2+\ldots+k^{2}=\frac{k(k+1)(2k+1)}6$.Find the smallest positive integer $k$ such that $1^2+2^2+3^2+\ldots+k^2$ is a multiple of $200$.


 

	\item (\textit{AIME II 2002}) Find the least positive integer $k$ for which the equation $\left\lfloor\frac{2002}{n}\right\rfloor=k$ has no integer solutions for $n$. (The notation $\lfloor x\rfloor$ means the greatest integer less than or equal to $x$.)


 

	\item (\textit{AIME II 2002}) Let $\mathcal{S}$ be the set $\lbrace1,2,3,\ldots,10\rbrace$. Let $n$ be the number of sets of two non-empty disjoint subsets of $\mathcal{S}$. (Disjoint sets are defined as sets that have no common elements.) Find the remainder obtained when $n$ is divided by $1000$.


 

	\item (\textit{AIME II 2002}) While finding the sine of a certain angle, an absent-minded professor failed to notice that his calculator was not in the correct angular mode. He was lucky to get the right answer. The two least positive real values of $x$ for which the sine of $x$ degrees is the same as the sine of $x$ radians are $\frac{m\pi}{n-\pi}$ and $\frac{p\pi}{q+\pi}$, where $m$, $n$, $p$, and $q$ are positive integers. Find $m+n+p+q$.


 

	\item (\textit{AIME II 2002}) Two distinct, real, infinite geometric series each have a sum of $1$ and have the same second term. The third term of one of the series is $1/8$, and the second term of both series can be written in the form $\frac{\sqrt{m}-n}p$, where $m$, $n$, and $p$ are positive integers and $m$ is not divisible by the square of any prime. Find $100m+10n+p$.


 

	\item (\textit{AIME II 2002}) A basketball player has a constant probability of $.4$ of making any given shot, independent of previous shots. Let $a_n$ be the ratio of shots made to shots attempted after $n$ shots. The probability that $a_{10}=.4$ and $a_n\le.4$ for all $n$ such that $1\le n\le9$ is given to be $p^aq^br/\left(s^c\right)$ where $p$, $q$, $r$, and $s$ are primes, and $a$, $b$, and $c$ are positive integers. Find $\left(p+q+r+s\right)\left(a+b+c\right)$.


 

	\item (\textit{AIME II 2002}) In triangle $ABC$, point $D$ is on $\overline{BC}$ with $CD=2$ and $DB=5$, point $E$ is on $\overline{AC}$ with $CE=1$ and $EA=3$, $AB=8$, and $\overline{AD}$ and $\overline{BE}$ intersect at $P$. Points $Q$ and $R$ lie on $\overline{AB}$ so that $\overline{PQ}$ is parallel to $\overline{CA}$ and $\overline{PR}$ is parallel to $\overline{CB}$. It is given that the ratio of the area of triangle $PQR$ to the area of triangle $ABC$ is $m/n$, where $m$ and $n$ are relatively prime positive integers. Find $m+n$.


 

	\item (\textit{AIME II 2002}) The perimeter of triangle $APM$ is $152$, and the angle $PAM$ is a right angle. A circle of radius $19$ with center $O$ on $\overline{AP}$ is drawn so that it is tangent to $\overline{AM}$ and $\overline{PM}$. Given that $OP=m/n$ where $m$ and $n$ are relatively prime positive integers, find $m+n$.


 

	\item (\textit{AIME II 2002}) Circles $\mathcal{C}_{1}$ and $\mathcal{C}_{2}$ intersect at two points, one of which is $(9,6)$, and the product of the radii is $68$. The x-axis and the line $y = mx$, where $m > 0$, are tangent to both circles. It is given that $m$ can be written in the form $a\sqrt {b}/c$, where $a$, $b$, and $c$ are positive integers, $b$ is not divisible by the square of any prime, and $a$ and $c$ are relatively prime. Find $a + b + c$.


 

\end{enumerate}